\chapter{Differential forms}
\label{chap:forms}

Differential forms on manifolds modeled on locally convex spaces and the exterior derivative,
following \cite[Section~3.6]{GN26}. Two features of the finite-dimensional theory are missing:
there is no coordinate description of a form by components in a chart, and there are no
partitions of unity, even on Banach manifolds. The definitions below only use local
representatives, and all statements are local. No completeness of the model space or of the
space of values is assumed.

Glöckner and Neeb define the exterior derivative through the Lie algebra differential of vector
fields \cite[Proposition~3.6.7, formula~(3.29)]{GN26} and show that it is given in charts by the
formula~(3.33) below. We take~(3.33) as the definition and check directly that it does not depend
on the chart: the terms with second derivatives of the transition maps cancel because these are
symmetric. The same kind of computation proves $\dd \circ \dd = 0$, as indicated in
\cite[Remark~3.6.10]{GN26}.

\section{Local forms}

Let $E$ and $F$ be Hausdorff locally convex spaces and $U \subseteq E$ open.

\begin{definition}[Local forms]
  \label{def:local-form}
  \lean{LocallyConvex.Manifolds.IsLocalForm}
  \leanok
  A \emph{local $F$-valued $n$-form} on $U$ is a map $\omega \colon U \times E^n \to F$ which is
  smooth in the sense of Michal--Bastiani and alternating $n$-linear in the last $n$ variables.
\end{definition}

\begin{definition}[Local exterior derivative]
  \label{def:local-d}
  \lean{LocallyConvex.Manifolds.dLoc}
  \leanok
  \uses{def:local-form}
  For a local $n$-form $\omega$,
  \[
    \dd\omega(y)(v_0, \dots, v_n) =
      \sum_{i=0}^{n} (-1)^i \, \partial_{v_i} \omega(y)(v_0, \dots, \widehat{v_i}, \dots, v_n),
  \]
  where $\partial_{v_i}$ is the directional derivative in the base point $y$
  \cite[(3.33)]{GN26}.
\end{definition}

\begin{proposition}
  \label{prop:local-d-form}
  \lean{LocallyConvex.Manifolds.IsLocalForm.dLoc, LocallyConvex.Manifolds.IsLocalForm.exists_partial}
  \leanok
  \uses{def:local-d}
  The exterior derivative of a local $n$-form is a local $(n+1)$-form.
\end{proposition}

\begin{proof}
  \leanok
  Smoothness: each term is the derivative $\dd\omega$ composed with a continuous linear map.
  For fixed $y$, the partial derivative $(u, w) \mapsto \partial_u \omega(y)(w)$ is linear in $u$
  and alternating $n$-linear in $w$, because $\omega(z)$ is so for every $z$ near $y$; then
  $\dd\omega(y)$ is its alternation, which is alternating $(n+1)$-linear.
\end{proof}

\begin{theorem}[$\dd \circ \dd = 0$, locally]
  \label{thm:local-dd}
  \lean{LocallyConvex.Manifolds.IsLocalForm.dLoc_dLoc,
    LocallyConvex.Manifolds.sum_neg_one_pow_smul_sum_eq_zero_of_symm}
  \leanok
  \uses{prop:local-d-form}
  For a local form $\omega$ on $U$, $\dd(\dd\omega) = 0$ on $U$.
\end{theorem}

\begin{proof}
  \leanok
  $\dd(\dd\omega)(y)(v_0, \dots, v_{n+1})$ is a double alternating sum of second derivatives
  $\partial_{v_i}\partial_{v_j}\omega(y)$, which are symmetric in $i, j$ (Schwarz's theorem for
  Michal--Bastiani maps). The terms cancel in pairs \cite[Remark~3.6.10]{GN26}.
\end{proof}

\begin{definition}[Local pullback]
  \label{def:local-pullback}
  \lean{LocallyConvex.Manifolds.pullbackLoc, LocallyConvex.Manifolds.IsLocalForm.pullback}
  \leanok
  \uses{def:local-form}
  For a smooth map $\tau \colon U \to V$ between open subsets of locally convex spaces and a
  local form $\omega$ on $V$,
  $(\tau^*\omega)(y)(v_1, \dots, v_n) = \omega(\tau y)(\dd\tau(y) v_1, \dots, \dd\tau(y) v_n)$.
  It is a local form on $U$.
\end{definition}

\begin{theorem}[Naturality, locally]
  \label{thm:local-naturality}
  \lean{LocallyConvex.Manifolds.IsLocalForm.dLoc_pullbackLoc,
    LocallyConvex.Manifolds.IsLocalForm.dMB_pullbackLoc,
    AlternatingMap.sum_neg_one_pow_smul_sum_update_eq_zero}
  \leanok
  \uses{def:local-pullback, def:local-d}
  $\dd(\tau^*\omega) = \tau^*(\dd\omega)$ on $U$.
\end{theorem}

\begin{proof}
  \leanok
  By the chain rule, $\partial_u (\tau^*\omega)(y)(w)$ is
  $\partial_{\dd\tau(y) u}\omega(\tau y)(\dd\tau(y) w)$ plus the sum over $k$ of
  $\omega(\tau y)(\dd\tau(y) w_1, \dots, \dd^2\tau(y)(w_k, u), \dots, \dd\tau(y) w_n)$; the second
  part uses that $\omega(\tau y)$ is multilinear. In the alternating sum, the first terms give
  $\tau^*(\dd\omega)$, and the terms with $\dd^2\tau(y)$ cancel because $\dd^2\tau(y)$ is
  symmetric.
\end{proof}

\section{Forms on manifolds}

Let $M$ be a smooth manifold modeled on $E$. A tangent vector at $x$ is given by its coordinates
in the preferred chart at $x$ (Definition~\ref{def:tangent-bundle}).

\begin{definition}[Differential forms]
  \label{def:diff-form}
  \lean{LocallyConvex.Manifolds.DifferentialForm, LocallyConvex.Manifolds.formLocalRep,
    LocallyConvex.Manifolds.DifferentialForm.localRep}
  \leanok
  \uses{def:tangent-bundle, def:local-form}
  An \emph{$F$-valued $n$-form} $\omega$ on $M$ assigns to each $x \in M$ an alternating
  $n$-linear map $\omega_x \colon (T_xM)^n \to F$ such that, for every chart $\varphi$ of the
  atlas, the local representative
  $(y, v_1, \dots, v_n) \mapsto \omega_{\varphi^{-1} y}(T\varphi^{-1} v_1, \dots,
  T\varphi^{-1} v_n)$ is smooth on $\varphi(U) \times E^n$ \cite[Definition~3.6.1]{GN26}. We
  write $\Omega^n(M, F)$ for the space of these forms.
\end{definition}

\begin{lemma}
  \label{lem:form-localrep}
  \lean{LocallyConvex.Manifolds.DifferentialForm.isLocalForm_localRep,
    LocallyConvex.Manifolds.DifferentialForm.localRep_eq_pullbackLoc,
    LocallyConvex.Manifolds.DifferentialForm.localRep_prefChart}
  \leanok
  \uses{def:diff-form, def:local-pullback}
  The local representatives of a form are local forms. On chart overlaps they are related by
  pullback along the transition map: $\omega_i = \tau^*\omega_j$ for $\tau = j \circ i^{-1}$.
\end{lemma}

\begin{proof}
  \leanok
  The coordinate changes of tangent vectors are the derivatives of the transition maps, and they
  compose by the chain rule.
\end{proof}

\begin{definition}[Exterior derivative]
  \label{def:diff-form-d}
  \lean{LocallyConvex.Manifolds.DifferentialForm.d, LocallyConvex.Manifolds.DifferentialForm.d_apply}
  \leanok
  \uses{def:diff-form, def:local-d}
  For $\omega \in \Omega^n(M, F)$, $(\dd\omega)_x$ is the local exterior derivative of the local
  representative of $\omega$ in the preferred chart at $x$, evaluated at the image of $x$.
\end{definition}

\begin{theorem}
  \label{thm:diff-form-d}
  \lean{LocallyConvex.Manifolds.DifferentialForm.localRep_d,
    LocallyConvex.Manifolds.DifferentialForm.formLocalRep_dFun}
  \leanok
  \uses{def:diff-form-d, lem:form-localrep, thm:local-naturality, prop:local-d-form}
  In every chart of the atlas, the local representative of $\dd\omega$ is the local exterior
  derivative of the local representative of $\omega$. In particular, $\dd\omega$ is a smooth
  $(n+1)$-form \cite[Proposition~3.6.7]{GN26}.
\end{theorem}

\begin{proof}
  \leanok
  Near a point, $\omega_i = \tau^*\omega_j$ with $\tau$ the transition map from the chart $i$ to
  the preferred chart $j$ at the point; apply Theorem~\ref{thm:local-naturality}.
\end{proof}

\begin{theorem}[$\dd \circ \dd = 0$]
  \label{thm:diff-form-dd}
  \lean{LocallyConvex.Manifolds.DifferentialForm.d_d}
  \leanok
  \uses{thm:diff-form-d, thm:local-dd}
  For every $\omega \in \Omega^n(M, F)$, $\dd(\dd\omega) = 0$ \cite[Proposition~3.6.9]{GN26}.
\end{theorem}

\begin{proof}
  \leanok
  In the preferred chart at $x$, the local representative of $\dd\omega$ is $\dd$ of the local
  representative of $\omega$ (Theorem~\ref{thm:diff-form-d}); apply
  Theorem~\ref{thm:local-dd}.
\end{proof}

\begin{definition}[Functions, closed and exact forms]
  \label{def:closed-exact}
  \lean{LocallyConvex.Manifolds.DifferentialForm.ofFunction,
    LocallyConvex.Manifolds.DifferentialForm.IsClosed,
    LocallyConvex.Manifolds.DifferentialForm.IsExact,
    LocallyConvex.Manifolds.DifferentialForm.IsExact.isClosed}
  \leanok
  \uses{def:diff-form-d, thm:diff-form-dd}
  A smooth function $f \colon M \to F$ is an $F$-valued $0$-form. A form $\omega$ is
  \emph{closed} if $\dd\omega = 0$ and \emph{exact} if $\omega = \dd\eta$ for some form $\eta$
  \cite[Definition~3.6.11]{GN26}. By Theorem~\ref{thm:diff-form-dd}, exact forms are closed.
\end{definition}

\begin{proposition}
  \label{prop:d-function}
  \lean{LocallyConvex.Manifolds.DifferentialForm.d_ofFunction_apply}
  \leanok
  \uses{def:closed-exact}
  For a smooth function $f$, $(\dd f)_x(v) = \dd(f \circ \varphi^{-1})(\varphi x)\, v$, where
  $\varphi$ is the preferred chart at $x$: the exterior derivative of $f$ is its differential
  \cite[Remark~3.6.12(a)]{GN26}.
\end{proposition}

\begin{proof}
  \leanok
  For $n = 0$ the alternating sum of Definition~\ref{def:local-d} has a single term.
\end{proof}

\section{Pullbacks}

\begin{definition}[Pullback]
  \label{def:diff-form-pullback}
  \lean{LocallyConvex.Manifolds.DifferentialForm.pullback,
    LocallyConvex.Manifolds.DifferentialForm.formLocalRep_pullbackFun}
  \leanok
  \uses{def:diff-form, def:tangent-map, def:local-pullback}
  For a smooth map $\varphi \colon N \to M$ and $\omega \in \Omega^n(M, F)$,
  $(\varphi^*\omega)_x(v_1, \dots, v_n) = \omega_{\varphi(x)}(T_x\varphi\, v_1, \dots,
  T_x\varphi\, v_n)$ \cite[Definition~3.6.6]{GN26}. In charts, the local representative of
  $\varphi^*\omega$ is the local pullback of the local representative of $\omega$ along the local
  representative of $\varphi$; in particular $\varphi^*\omega$ is a smooth form.
\end{definition}

\begin{theorem}[Naturality of $\dd$]
  \label{thm:diff-form-naturality}
  \lean{LocallyConvex.Manifolds.DifferentialForm.d_pullback}
  \leanok
  \uses{def:diff-form-pullback, thm:diff-form-d, thm:local-naturality}
  For a smooth map $\varphi \colon N \to M$ and $\omega \in \Omega^n(M, F)$,
  $\dd(\varphi^*\omega) = \varphi^*(\dd\omega)$ \cite[Lemma~3.6.13]{GN26}.
\end{theorem}

\begin{proof}
  \leanok
  In the preferred charts at $x$ and $\varphi(x)$, this is Theorem~\ref{thm:local-naturality}
  for the local representative of $\varphi$.
\end{proof}
